\documentclass[10pt,a4paper]{amsart}
\usepackage[margin=19mm]{geometry}
\usepackage{amsmath,amssymb,mathtools}
\usepackage[scr=rsfs,cal=euler]{mathalfa}
\usepackage{xfrac}
\usepackage[unicode,hidelinks,bookmarks=false]{hyperref}
\newcommand{\ie}{i.e.\ }
\newcommand{\cf}{cf.\ }
\newcommand{\resp}{respectively\ }
\newcommand{\Subsection}{\subsection}
\providecommand{\nobreakdash}{\nobreak\hbox{-}}
\setlength{\emergencystretch}{2em}
\multlinegap0pt
\usepackage[shortlabels]{enumitem}
\usepackage{tikz}
\usepackage{amssymb}
\usepackage{xfrac}
\usepackage{dsfont}
\usepackage{mathtools}
\newtheorem{proposition}{Proposition}
\newtheorem{lemma}{Lemma}
\DeclareMathOperator{\Jac}{Jac}
\newcommand{\equivk}{\overset{\mathcal{K}}{\sim}}
\newcommand{\equivr}{\overset{\mathcal{R}}{\sim}}
\DeclareMathOperator{\fpt}{fpt}
\newcommand{\kx}{k[[\underline{x}]]}
\DeclareMathOperator{\ord}{ord}
\title{Bounds on the $F$-Pure Threshold of Isolated Hypersurface Singularities}
\author{Yotam Svoray}
\date{Source: arXiv:2606.04014v2 (CC BY 4.0)}
\begin{document}
\maketitle

\noindent\emph{Source and licence.} Bounds on the $F$-Pure Threshold of Isolated Hypersurface Singularities. Version arXiv:2606.04014v2 (CC BY 4.0), distributed by arXiv under the Creative Commons Attribution 4.0 International licence, http://creativecommons.org/licenses/by/4.0/, as recorded in the arXiv licence metadata for this exact version and shown on the abstract page https://arxiv.org/abs/2606.04014v2. Full licence notice: \texttt{licenses/arxiv-2606.04014-CC-BY-4.0.txt}.\par\smallskip

\noindent\textit{Editorial scope.} Counted unit: the article's lemma lem:c(f,Jac)\_fpt (source0.tex lines 232--238). The article's complete original proof of it is delivered with it (source0.tex lines 240--260, one \texttt{proof} environment of 21 lines). No mathematics is written, repaired or re-derived: the statement and the proof are the article's own lines, in the article's own notation, and no retained line is substituted or re-flowed.  Retained uncounted as the source-local context the proof needs: lines 156--167.  Nothing was omitted from any retained span, so the package carries no omission record.  No retained line contains a citation, so the wrapper carries no bibliography and no bibliography entry.  No retained line contains a figure environment, an image-inclusion command or drawing code, so no image is delivered and the package declares no asset dependency.  Editorial formatting only, in addition: an AMS \texttt{amsart} preamble that restates the article's own declarations and macros used by the retained lines, namely the environments \texttt{proposition}, \texttt{lemma} on separate counters, together with the standard packages those lines need.  The retained environments print in this document as Proposition 1, Lemma 1 in order of appearance, whereas the article prints the same items with the numbering of its own full document.  The complete native source of the article as distributed in the arXiv e-print for this exact version is bundled unchanged as \texttt{sources/202019/source0.tex}.
\begin{proposition}\label{prop:mt_basic}
Let $f\in\mathfrak{m}\subset\kx$. Then:
\begin{enumerate}
    \item If $\tau(f)<\infty$, then $\mu(f)<\infty$ if and only if $f\in\sqrt{\Jac(f)}$. In addition, if $f^N\in\Jac(f)$, then $\tau(f)\leq\mu(f)\leq N\cdot\tau(f)$.
    \item (Detecting Smoothness) $\mu(f)=0$ if and only if $\tau(f)=0$ if and only if $f\equivr x_1$.
    \item (Stability under $\equivk$) Given $g\in\kx$, if $f\equivk g$, then $\tau(f)=\tau(g)$.
    \item (Stability under $\equivr$) Given $g\in\kx$, if $f\equivr g$, then $\mu(f)=\mu(g)$.
    \item (Contact Determinacy) $\tau(f)<\infty$ if and only if there exists $N>0$ such that for every $g\in\kx$ with $g-f\in\mathfrak{m}^{N+1}$ we have $f\equivk g$. In this case we can choose $N=2\tau(f)-\ord(f)+2$.
    \item (Right Determinacy) $\mu(f)<\infty$ if and only if there exists $N>0$ such that for every $g\in\kx$ with $g-f\in\mathfrak{m}^{N+1}$ we have $f\equivr g$. In this case we can choose $N=2\mu(f)-\ord(f)+2$.
    \item If $\mu(f)<\infty$ (resp.\ $\tau(f)<\infty$), then $\mathfrak{m}^{\mu(f)}\subset\Jac(f)$ (resp.\ $\mathfrak{m}^{\tau(f)}\subset\Jac(f)+\langle f\rangle$).
\end{enumerate}
\end{proposition}

\begin{lemma}\label{lem:c(f,Jac)_fpt}
Let $0\neq f\in\mathfrak{m}$. Then:
\begin{enumerate}
    \item If $\mu(f)<\infty$, then $\fpt(f)\mu(f)\geq c(f,\Jac(f))$.
    \item If $0<\tau(f)<\infty$, then $\fpt(f)\tau(f)\geq c(f,\Jac(f)+\langle f\rangle)$.
\end{enumerate}
\end{lemma}

\begin{proof}
If $\mu(f)<\infty$, Proposition~\ref{prop:mt_basic} gives $\mathfrak{m}^{\mu(f)}\subset\Jac(f)$, and hence
\[
(\mathfrak{m}^{[p^e]})^{\mu(f)}\subset\Jac(f)^{[p^e]}.
\]
Since $f^{\nu_e(f)+1}\in\mathfrak{m}^{[p^e]}$, it follows that
\[
\mu(f)(\nu_e(f)+1)\geq\nu_e(f,\Jac(f))+1.
\]
Dividing by $p^e$ and taking the limit proves the first assertion.  Similarly, if $0<\tau(f)<\infty$, then $\mathfrak{m}^{\tau(f)}\subset\Jac(f)+\langle f\rangle$, so
\[
(\mathfrak{m}^{[p^e]})^{\tau(f)}
\subset(\Jac(f)+\langle f\rangle)^{[p^e]}.
\]
The same argument gives
\[
\tau(f)(\nu_e(f)+1)
\geq\nu_e(f,\Jac(f)+\langle f\rangle)+1,
\]
and the second assertion follows after dividing by $p^e$ and taking the limit.
\end{proof}

\end{document}
